\chapter*{Resumen}

Este Trabajo Fin de Grado presenta el diseño e implementación de \textit{NeuroForge}, un videojuego de estrategia por turnos con mecánicas de información oculta desarrollado sobre el motor Godot. El videojuego está inspirado en el clásico juego de mesa \textit{Stratego}, adaptando sus mecánicas fundamentales a un entorno digital con ambientación de ciencia ficción futurista.

El proyecto abarca el ciclo de vida completo del desarrollo de software, desde el análisis de requisitos y el diseño de la arquitectura hasta la implementación, las pruebas y la documentación. La arquitectura adoptada separa en tres capas la lógica del videojuego, la interfaz de usuario y la infraestructura, lo que ha facilitado el mantenimiento y la escalabilidad del sistema a lo largo de las cinco iteraciones de desarrollo. Se han aplicado patrones de diseño orientados a objetos como \textit{Singleton}, \textit{Strategy}, \textit{Fachada} y \textit{Template Method} para garantizar una estructura robusta y desacoplada.

Como componente avanzado, se ha desarrollado un agente de inteligencia artificial entrenado mediante aprendizaje por refuerzo, utilizando la biblioteca Gymnasium como interfaz del entorno y el algoritmo MaskablePPO como método de entrenamiento. El agente aprende a jugar mediante una estrategia de \textit{self-play}, enfrentándose progresivamente a versiones anteriores de sí mismo, lo que le permite desarrollar comportamientos estratégicos sin necesidad de una función de evaluación diseñada  manualmente.

El resultado es un videojuego completo y jugable, con una interfaz visual pulida, efectos sonoros, múltiples menús y un oponente controlado por inteligencia artificial, distribuible como ejecutable autocontenido para Windows.

\vspace{0.5cm}
\noindent \textbf{Palabras clave:} videojuego, estrategia por turnos, información oculta, Godot, aprendizaje por refuerzo, Gymnasium, MaskablePPO, \textit{self-play}, arquitectura en capas, patrones de diseño.






